\documentclass{article}
% Retain the A3 composition as the logical typesetting canvas; the physical
% output below is exactly 250 x 350 mm. Fonts below account for the 5/6 scale.
\usepackage[a3paper,portrait,hmargin=6mm,top=7mm,bottom=7mm]{geometry}
\usepackage{pgfpages}
% Uniform 5/6 scaling: 297 x 420 mm becomes 247.5 x 350 mm.
% Centering adds 1.25 mm at each side to absorb the aspect-ratio difference.
\pgfpagesdeclarelayout{poster-board-eighth}{}{%
  \pgfpagesphysicalpageoptions{logical pages=1,
    physical width=250mm,physical height=350mm}%
  \pgfpageslogicalpageoptions{1}{scale=0.833333333,
    center=\pgfpoint{125mm}{175mm}}%
}
\pgfpagesuselayout{poster-board-eighth}
\usepackage[T1]{fontenc}
\usepackage{lmodern}
\usepackage{amsmath,amssymb}
% Use scalable mathematical operators at the reference text sizes.
\input{omxlmex.fd}
\DeclareFontShape{OMX}{lmex}{m}{n}{<->lmex10}{}
\usepackage{xcolor,colortbl}
\usepackage{tikz}
\usepackage{tabularx,array,collcell}
% Text tables use a symmetric minimum row strut around the cell centers.
% Keep its total height and leave mathematical arrays' internal spacing intact.
\makeatletter
\pretocmd{\@mkpream}{%
  \ifx\@acol\@tabacol
    \@tempdima=\dimexpr(\ht\@arstrutbox+\dp\@arstrutbox)/2\relax
    \setbox\@arstrutbox=\hbox{\vrule width0pt height\@tempdima depth\@tempdima}%
  \fi
  }{}{\PackageError{notation-posters}{Cannot center table struts}{}}
\makeatother
\pagestyle{empty}
% Central group palettes: dark ink, instructional accent, and quiet table rules.
% Group A retains the original colors; group B uses coordinated cool teals.
\definecolor{PosterAInk}{HTML}{172C3C}
\definecolor{PosterAAccent}{HTML}{BF604D}
\definecolor{PosterARule}{HTML}{A5AFB6}
\definecolor{PosterBInk}{HTML}{17383B}
\definecolor{PosterBAccent}{HTML}{287C7A}
\definecolor{PosterBRule}{HTML}{91AAA9}
% Existing semantic aliases also drive the unchanged panel tints and markers.
% Select after a page break so the preceding page retains its own palette.
\newcommand{\PosterGroup}[1]{%
  \colorlet{navy}{Poster#1Ink}%
  \colorlet{coral}{Poster#1Accent}%
  \colorlet{rulecolor}{Poster#1Rule}%
}
\PosterGroup{A}
% One restrained marker for the specific notation introduced in each row.
\newif\ifGTenHighlights
\GTenHighlightstrue
% \GTenHighlightsfalse % Uncomment to disable instructional highlights.
% Rounded highlights keep even padding and the current math style.
% Native symbol metrics preserve spacing and baselines; padding uses surrounding clearance.
\newcommand{\sequenceindex}[1]{\ifGTenHighlights\mathpalette\sequenceindexbox{#1}\fi#1}
\newcommand{\sequenceindexbox}[2]{\begingroup
  % Paint only the background; the original symbol keeps its native math atom,
  % italic correction, script attachment, and dimensions outside this overlay.
  \raisebox{0pt}[0pt][0pt]{%
  \makebox[0pt][l]{\hspace{-1.2pt}\tikz[baseline=(sequenceindex.base)]\node[
    fill=coral!12,fill opacity=0.65,draw=none,rounded corners=2.6pt,
    inner xsep=1.2pt,inner ysep=2.2pt,outer sep=0pt,
    minimum width=0pt,minimum height=0pt
  ] (sequenceindex) {$\mathsurround=0pt #1\phantom{#2}$};}%
  }%
  \endgroup}
% Keep integral bounds on the native operator, including its italic correction;
% the shared renderer paints behind the integral sign only.
\newcommand{\notationintegral}[1]{\mathop{%
  \mathopen{\ifGTenHighlights\mathpalette\sequenceindexbox{\int}\fi}#1}\nolimits}
% Highlight the fraction rule alone, using the same renderer and toggle.
% Keep the native fraction (and its numerator/denominator spacing) intact.
\newsavebox{\FractionMeasure}
\newlength{\FractionBarWidth}
\newlength{\FractionAxis}
\newlength{\FractionRule}
\newcommand{\notationfraction}[2]{\mathinner{%
  \mathopen{\ifGTenHighlights\mathpalette\notationfractionbar{{#1}{#2}}\fi}%
  \frac{#1}{#2}}}
\newcommand{\notationfractionbar}[2]{\notationfractionmeasure{#1}#2}
\newcommand{\notationfractionmeasure}[3]{\begingroup
  \sbox{\FractionMeasure}{$\mathsurround=0pt #1\frac{#2}{#3}$}%
  \setlength{\FractionBarWidth}{\dimexpr\wd\FractionMeasure-2\nulldelimiterspace\relax}%
  \ifx#1\displaystyle
    \setlength{\FractionAxis}{\fontdimen22\textfont2}%
    \setlength{\FractionRule}{\fontdimen8\textfont3}%
  \else\ifx#1\textstyle
    \setlength{\FractionAxis}{\fontdimen22\textfont2}%
    \setlength{\FractionRule}{\fontdimen8\textfont3}%
  \else\ifx#1\scriptstyle
    \setlength{\FractionAxis}{\fontdimen22\scriptfont2}%
    \setlength{\FractionRule}{\fontdimen8\scriptfont3}%
  \else
    \setlength{\FractionAxis}{\fontdimen22\scriptscriptfont2}%
    \setlength{\FractionRule}{\fontdimen8\scriptscriptfont3}%
  \fi\fi\fi
  \raisebox{\dimexpr\FractionAxis-\FractionRule/2\relax}[0pt][0pt]{%
    \makebox[0pt][l]{\hspace{\nulldelimiterspace}%
      \sequenceindexbox{#1}{\rule{\FractionBarWidth}{\FractionRule}}}}%
  \endgroup}

\color{navy}
\setlength{\parindent}{0pt}
\setlength{\parskip}{0pt}
\setlength{\arrayrulewidth}{0.3pt}
% Central typography controls, explicitly applied in table cells.
% Printed sizes are 5/6 of these logical sizes: body 10.83 pt,
% examples 10.67 pt, symbols 11.67 pt, table labels/notes 10.33 pt.
\newcommand{\TitleFont}{\fontsize{27}{30}\selectfont\bfseries}
\newcommand{\KickerFont}{\fontsize{13}{15}\selectfont\bfseries}
\newcommand{\HeadingFont}{\fontsize{17}{20}\selectfont\bfseries}
\newcommand{\DefinitionFont}{\fontsize{13}{15}\selectfont}
\newcommand{\SymbolFont}{\fontsize{14}{16}\selectfont}
\newcommand{\ExampleFont}{\fontsize{12.8}{14.8}\selectfont}
\newcommand{\TableHeadFont}{\fontsize{12.4}{14.4}\selectfont\bfseries}
\newcommand{\FooterFont}{\fontsize{12.4}{14.4}\selectfont}
\DeclareMathSizes{12.6}{12.6}{8.8}{6.6}
\DeclareMathSizes{14}{14}{9.8}{7}
\DeclareMathSizes{13}{13}{9.1}{6.8}
\DeclareMathSizes{12.8}{12.8}{9}{6.8}

\newlength{\SymbolWidth}
\newlength{\MeaningWidth}
\newlength{\RowPadding}
\setlength{\RowPadding}{0.7mm}
% Measure the complete cell (including wrapped text and tall mathematics),
% then reserve the same clearance above and below its actual bounds.
\newsavebox{\CellContent}
% All table helpers share the same vertical center and local padding.
% An extra 0.4 mm on each side gives A3 rows consistent breathing room.
\newcommand{\paddedcell}[1]{\centeredcell{#1}}
\newcommand{\centeredtextcolumn}[1]{p{#1}}
\newcommand{\contextcell}[1]{%
  \sbox{\CellContent}{#1}%
  \raisebox{\dimexpr(\dp\CellContent-\ht\CellContent)/2\relax}{\usebox{\CellContent}}}
\newcommand{\centeredcell}[1]{%
  \sbox{\CellContent}{\begin{minipage}[t]{\linewidth}\raggedright #1\end{minipage}}%
  \raisebox{\dimexpr(\dp\CellContent-\ht\CellContent)/2\relax}%
    [\dimexpr(\ht\CellContent+\dp\CellContent)/2+\RowPadding+0.4mm\relax]%
    [\dimexpr(\ht\CellContent+\dp\CellContent)/2+\RowPadding+0.4mm\relax]%
    {\usebox{\CellContent}}%
}
% Text-style arrays align relations without enlarging fractions or operators.
% Each r/E pair holds a left-hand side and its relation/right-hand side.
% E supplies an empty math atom to retain normal spacing before a relation.
\newcommand{\ExampleLineStretch}{1.12}
% Keep each concept and its explanation in one measured Meaning cell.
% Paragraphs work in the cell minipages and the paired/triple parboxes alike.
\newcommand{\meaningexplanation}[1]{\textcolor{gray}{\textit{#1}}}
\newcommand{\conceptmeaning}[2]{\begingroup
  % Keep the concept and gray explanation close with readable A3 leading.
  % Inherit the current body leading, including the dense calculus page.
  #1\par #2\rule[-1pt]{0pt}{1pt}\par
  \endgroup}
\newsavebox{\ExampleContent}
% A second spacing pass is enabled only in the compact tables.
\newcommand{\TightExampleSpacing}{}
\newcommand{\exmath}[2][r@{}E]{\begingroup
  \renewcommand{\arraystretch}{\ExampleLineStretch}%
  \ifcsname NC@find@E\endcsname\else\newcolumntype{E}{>{{}}l}\fi%
  \sbox{\ExampleContent}{$\begin{array}[t]{@{}#1@{}}#2\end{array}$}%
  % Only overwide examples receive a small math-spacing correction.
  % This preserves the font size, baselines, and all fitting examples.
  \ifdim\wd\ExampleContent>\linewidth
    \thickmuskip=4mu\relax\medmuskip=3mu\relax
    \sbox{\ExampleContent}{$\begin{array}[t]{@{}#1@{}}#2\end{array}$}%
  \fi
  \ifdim\wd\ExampleContent>\linewidth
    \TightExampleSpacing
    \sbox{\ExampleContent}{$\begin{array}[t]{@{}#1@{}}#2\end{array}$}%
  \fi
  \usebox{\ExampleContent}\endgroup}
\newcommand{\symrow}[3]{%
  \centeredcell{\centering\SymbolFont\ensuremath{\displaystyle #1}}&%
  \centeredcell{\DefinitionFont #2}&\centeredcell{\ExampleFont #3}\\\hline
}
% Composite rows retain one shared example array for horizontal relation alignment.
% Measure each logical row separately; center all its cells on that row's axis.
\newsavebox{\PairFirst}
\newsavebox{\PairSecond}
\newsavebox{\PairThird}
\newsavebox{\ExampleFirst}
\newsavebox{\ExampleSecond}
\newsavebox{\ExampleThird}
\newsavebox{\CurrentExample}
\newsavebox{\CurrentRow}
\newlength{\ExampleWidth}
\newcommand{\MeasuredLineStretch}{1}
\newcommand{\subrowcell}[1]{\begingroup
  \setlength{\RowPadding}{0pt}\centeredcell{#1}\endgroup}
\newcommand{\pairstrut}[1]{\rule[-\dp#1]{0pt}{\dimexpr\ht#1+\dp#1\relax}}
\newcommand{\measurelogicalrow}[6]{%
  \global\setbox#1=\hbox{\ExampleFont\exmath[#3]{#6}}%
  \global\setbox#2=\hbox{%
    {\setlength{\linewidth}{\SymbolWidth}\subrowcell{\centering\SymbolFont$\displaystyle #4$}}%
    {\setlength{\linewidth}{\MeaningWidth}\subrowcell{\DefinitionFont #5}}%
    \raisebox{\dimexpr(\dp#1-\ht#1)/2\relax}{\usebox{#1}}}%
}
\newcommand{\pairlines}[3]{%
  \renewcommand{\arraystretch}{0}%
  \begin{tabular}[t]{@{}#1@{}}%
  \pairstrut{\PairFirst}\subrowcell{#2}\\%
  \pairstrut{\PairSecond}\subrowcell{#3}%
  \end{tabular}%
}
\newcommand{\triplelines}[4]{%
  \renewcommand{\arraystretch}{0}%
  \begin{tabular}[t]{@{}#1@{}}%
  \pairstrut{\PairFirst}\subrowcell{#2}\\%
  \pairstrut{\PairSecond}\subrowcell{#3}\\%
  \pairstrut{\PairThird}\subrowcell{#4}%
  \end{tabular}%
}
% Apply one translation to every math column in a logical row, preserving
% the relative baselines of its relations, fractions, and other mathematics.
\newcommand{\compositemathcell}[1]{%
  \pairstrut{\CurrentRow}%
  \raisebox{\dimexpr(\dp\CurrentExample-\ht\CurrentExample)/2\relax}{\renewcommand{\arraystretch}{\MeasuredLineStretch}$#1$}}
\newcommand{\compositerelationcell}[1]{\compositemathcell{{}#1}}
\newcolumntype{R}{>{\collectcell\compositemathcell}r<{\endcollectcell}}
\newcolumntype{L}{>{\collectcell\compositemathcell}l<{\endcollectcell}}
\newcolumntype{Q}{>{\collectcell\compositerelationcell}l<{\endcollectcell}}
\ExplSyntaxOn
\newcommand{\compositecolumns}[1]{
  \tl_set:Nn \l_tmpa_tl {#1}
  \tl_replace_all:Nnn \l_tmpa_tl {r}{R}
  \tl_replace_all:Nnn \l_tmpa_tl {l}{L}
  \tl_replace_all:Nnn \l_tmpa_tl {E}{Q}
  \cs_set:Npx \CompositeColumns {\exp_not:V \l_tmpa_tl}
}
\ExplSyntaxOff
\newcommand{\selectlogicalrow}[2]{%
  \global\setbox\CurrentRow=\copy#1%
  \global\setbox\CurrentExample=\copy#2}
\newcommand{\preparecomposite}{%
  \renewcommand{\ExampleLineStretch}{\MeasuredLineStretch}%
  \setlength{\ExampleWidth}{\dimexpr\linewidth-\SymbolWidth-\MeaningWidth-4\tabcolsep\relax}%
  \setlength{\linewidth}{\ExampleWidth}}
\newcommand{\triplerow}[9]{%
  \noalign{\begingroup\preparecomposite
    \measurelogicalrow{\ExampleFirst}{\PairFirst}{r@{}E}{#1}{#4}{#7}%
    \measurelogicalrow{\ExampleSecond}{\PairSecond}{r@{}E}{#2}{#5}{#8}%
    \measurelogicalrow{\ExampleThird}{\PairThird}{r@{}E}{#3}{#6}{#9}%
  \endgroup}%
  \paddedcell{\centering\SymbolFont\triplelines{c}%
    {\centering$\displaystyle #1$}{\centering$\displaystyle #2$}{\centering$\displaystyle #3$}}&%
  \paddedcell{\DefinitionFont\triplelines{l}{#4}{#5}{#6}}&%
  \paddedcell{\ExampleFont\renewcommand{\ExampleLineStretch}{0}%
    \selectlogicalrow{\PairFirst}{\ExampleFirst}%
    \exmath[R@{}Q]{\noalign{\selectlogicalrow{\PairFirst}{\ExampleFirst}}#7\\
      \noalign{\selectlogicalrow{\PairSecond}{\ExampleSecond}}#8\\
      \noalign{\selectlogicalrow{\PairThird}{\ExampleThird}}#9}}\\\hline
}
\newcommand{\pairedrow}[6]{%
  \pairedrowwithcols{r@{}E@{\quad}r@{}E}{#1}{#2}{#3}{#4}{#5}{#6}}
\newcommand{\pairedrowwithcols}[7]{%
  \noalign{\begingroup\preparecomposite
    \measurelogicalrow{\ExampleFirst}{\PairFirst}{#1}{#2}{#4}{#6}%
    \measurelogicalrow{\ExampleSecond}{\PairSecond}{#1}{#3}{#5}{#7}%
  \endgroup}%
  \paddedcell{\centering\SymbolFont\pairlines{c}{\centering$\displaystyle #2$}{\centering$\displaystyle #3$}}&%
  \paddedcell{\DefinitionFont\pairlines{l}{#4}{#5}}&%
  \paddedcell{\ExampleFont\renewcommand{\ExampleLineStretch}{0}%
    \compositecolumns{#1}%
    \def\SingleExampleColumn{l}\def\TheseExampleColumns{#1}%
    \ifx\SingleExampleColumn\TheseExampleColumns
      \def\FirstExample{{#6}}\def\SecondExample{{#7}}%
    \else\def\FirstExample{#6}\def\SecondExample{#7}\fi%
    \selectlogicalrow{\PairFirst}{\ExampleFirst}%
    \expandafter\exmath\expandafter[\CompositeColumns]{\noalign{\selectlogicalrow{\PairFirst}{\ExampleFirst}}\FirstExample\\
      \noalign{\selectlogicalrow{\PairSecond}{\ExampleSecond}}\SecondExample}}\\\hline
}
% Compact contexts use the existing palette and the table's example type.
\newcommand{\ContextHeading}{Example contexts}
\newcommand{\examplecontexts}[1]{%
  \begingroup
  \setlength{\fboxsep}{0.5mm}%
  \setlength{\fboxrule}{0.35pt}%
  \noindent\fcolorbox{coral}{coral!10}{%
    \begin{minipage}{\dimexpr\linewidth-2\fboxsep-2\fboxrule\relax}
      \ExampleFont\color{navy}%
      {\centering\textbf{\ContextHeading}\par}
      #1\par
    \end{minipage}}\par
  \endgroup
  \vspace{0.7mm}%
}
% Optional contexts appear after the heading rule and before the table.
\newcommand{\SymbolColumnFraction}{.24}
\newcommand{\MeaningColumnFraction}{.32}
% A heading always spans the page; parallel tables belong to that section.
\newcommand{\sectionheading}[1]{%
  {\HeadingFont\color{coral}#1\par}%
  \vspace{0.4mm}{\color{navy}\hrule height 0.5pt}\vspace{0.8mm}%
}
\newcommand{\notationtable}[1]{%
  \begingroup
  \setlength{\SymbolWidth}{\SymbolColumnFraction\linewidth}%
  \setlength{\MeaningWidth}{\MeaningColumnFraction\linewidth}%
  \setlength{\tabcolsep}{0.5mm}%
  \renewcommand{\arraystretch}{1}%
  \arrayrulecolor{rulecolor}%
  \begin{tabularx}{\linewidth}{@{}>{\centering\arraybackslash\SymbolFont}p{\SymbolWidth}>{\raggedright\arraybackslash\DefinitionFont}p{\MeaningWidth}>{\raggedright\arraybackslash\ExampleFont}X@{}}
  \paddedcell{\centering\TableHeadFont Symbol}&\paddedcell{\TableHeadFont Meaning}&\paddedcell{\TableHeadFont Examples}\\\hline
  #1
  \end{tabularx}%
  \endgroup
}
% Read the complete left table, then the right table, before the next section.
\newcommand{\paralleltables}[2]{%
  \noindent\begin{minipage}[t]{.49\linewidth}\vspace{0pt}\notationtable{#1}\end{minipage}\hfill
  \begin{minipage}[t]{.49\linewidth}\vspace{0pt}\notationtable{#2}\end{minipage}\par
  \vspace{1mm}%
}
\newcommand{\sectionbox}[4][]{%
  \sectionheading{#2}%
  \if\relax\detokenize{#1}\relax\else\examplecontexts{#1}\fi
  \paralleltables{#3}{#4}%
}
% Dedicated two-column reference for letter names.
\newlength{\LetterMeaningWidth}
\newcommand{\letterrow}[2]{%
  \paddedcell{\centering\SymbolFont\ensuremath{#1}}&%
  \paddedcell{\DefinitionFont #2}\\\hline
}
\newcommand{\letterreference}[1]{%
  \begingroup
  \setlength{\tabcolsep}{0.7mm}%
  \setlength{\LetterMeaningWidth}{\dimexpr\linewidth-.29\linewidth-2\tabcolsep\relax}%
  \renewcommand{\arraystretch}{0.8}%
  \arrayrulecolor{rulecolor}%
  \begin{tabularx}{\linewidth}{@{}>{\centering\arraybackslash\SymbolFont}p{.29\linewidth}>{\raggedright\arraybackslash\DefinitionFont}X@{}}
  \paddedcell{\centering\TableHeadFont Symbol}&\paddedcell{\TableHeadFont Name / description}\\\hline
  #1
  \end{tabularx}%
  \endgroup\par\vspace{0.8mm}%
}

% Parallel tables retain their measured subrows, local widths, and padding.
\newcommand{\compacttables}{%
  \renewcommand{\TightExampleSpacing}{\thinmuskip=0mu\relax\medmuskip=1mu\relax\thickmuskip=1mu\relax}%
  \setlength{\RowPadding}{0.9mm}%
  \renewcommand{\MeasuredLineStretch}{1}%
  \renewcommand{\SymbolColumnFraction}{.22}%
  \renewcommand{\MeaningColumnFraction}{.35}%
  \renewcommand{\pairedrow}[6]{\pairedrowwithcols{r@{}E}{##1}{##2}{##3}{##4}{##5}{##6}}%
  \renewcommand{\symrow}[3]{%
    \paddedcell{\centering\SymbolFont\ensuremath{\displaystyle ##1}}&%
    \paddedcell{\DefinitionFont ##2}&\paddedcell{\ExampleFont ##3}\\\hline}%
}
% Letter composites share the same independent centers and outer padding.
\newcommand{\measureletterrow}[3]{%
  \global\setbox#1=\hbox{%
    {\setlength{\linewidth}{.29\linewidth}\subrowcell{\centering\SymbolFont$#2$}}%
    {\setlength{\linewidth}{\LetterMeaningWidth}\subrowcell{\DefinitionFont #3}}}}
\newcommand{\letterpair}[4]{%
  \noalign{\measureletterrow{\PairFirst}{#1}{#3}%
    \measureletterrow{\PairSecond}{#2}{#4}}%
  \paddedcell{\centering\SymbolFont\pairlines{c}{\centering$#1$}{\centering$#2$}}&%
  \paddedcell{\DefinitionFont\pairlines{l}{#3}{#4}}\\\hline}
\newcommand{\lettertriple}[6]{%
  \noalign{\measureletterrow{\PairFirst}{#1}{#4}%
    \measureletterrow{\PairSecond}{#2}{#5}%
    \measureletterrow{\PairThird}{#3}{#6}}%
  \paddedcell{\centering\SymbolFont\triplelines{c}{\centering$#1$}{\centering$#2$}{\centering$#3$}}&%
  \paddedcell{\DefinitionFont\triplelines{l}{#4}{#5}{#6}}\\\hline}

% Each poster has its own fixed page and repeated identity.
\newcommand{\posterheader}[1]{%
  {\TitleFont\color{coral}MATHEMATICS\hfill\KickerFont #1\par}%
  \vspace{2mm}{\color{navy}\hrule height 0.8pt}\vspace{3mm}%
  \color{navy}%
}
\begin{document}
\posterheader{A-1\quad FOUNDATIONS}
\begingroup
\begingroup
% General notation uses two independently headed tables in reading order.
\setlength{\RowPadding}{1.1mm}
\renewcommand{\MeasuredLineStretch}{1}
\renewcommand{\SymbolColumnFraction}{.26}
\renewcommand{\MeaningColumnFraction}{.33}
{\HeadingFont\color{coral}General notation\par}
\vspace{0.4mm}{\color{navy}\hrule height 0.5pt}\vspace{0.4mm}
\noindent\begin{minipage}[t]{.49\linewidth}
% Balance the shorter left table using cell clearance only.
\setlength{\RowPadding}{0.9mm}
\vspace{0pt}
\notationtable{
\triplerow{\mathrel{\sequenceindex{\mathord{=}}}}{\mathrel{\sequenceindex{\mathord{\neq}}}}{\mathrel{\sequenceindex{\mathord{\approx}}}}{Equal}{Not equal}{Approximately}
  {2+3&\mathrel{\sequenceindex{\mathord{=}}}5\quad 4-1\mathrel{\sequenceindex{\mathord{=}}}3}{2&\mathrel{\sequenceindex{\mathord{\neq}}}3\quad 5+1\mathrel{\sequenceindex{\mathord{\neq}}}7}{3.14&\mathrel{\sequenceindex{\mathord{\approx}}}\pi \quad\>\>\> \sqrt2\mathrel{\sequenceindex{\mathord{\approx}}}1.41}

\pairedrow{\mathrel{\sequenceindex{\mathord{<}}} \quad \mathrel{\sequenceindex{\mathord{>}}}}{\mathrel{\sequenceindex{\mathord{\nless}}} \quad \mathrel{\sequenceindex{\mathord{\ngtr}}}}
{Less | Greater than}
{Not Less | Greater}
{2&\mathrel{\sequenceindex{\mathord{<}}}3 & \qquad 5&\mathrel{\sequenceindex{\mathord{>}}}4}{2&\mathrel{\sequenceindex{\mathord{\nless}}}2 &\qquad  4&\mathrel{\sequenceindex{\mathord{\ngtr}}}5}

\pairedrow{\mathrel{\sequenceindex{\mathord{\leq}}} \quad \mathrel{\sequenceindex{\mathord{\geq}}}}{\mathrel{\sequenceindex{\mathord{\nleq}}} \quad \mathrel{\sequenceindex{\mathord{\ngeq}}}}
{Less | Greater than or equal to}
{Not Less | Greater than or equal to}
{2&\mathrel{\sequenceindex{\mathord{\leq}}}2 &\qquad  5&\mathrel{\sequenceindex{\mathord{\geq}}}4}{3&\mathrel{\sequenceindex{\mathord{\nleq}}}2 &\qquad  4&\mathrel{\sequenceindex{\mathord{\ngeq}}}5}

\symrow{\begin{gathered}ab=a(b)\\{}=a\mathbin{\sequenceindex{\cdot}} b=a\mathbin{\sequenceindex{\times}} b\end{gathered}}
{Multiplication: \\ \meaningexplanation{juxtaposition, parentheses, dot, or cross}}
{\exmath{3x&=3\mathbin{\sequenceindex{\cdot}} x=3\mathbin{\sequenceindex{\times}} x\\2(4)&=2\mathbin{\sequenceindex{\cdot}}4=2\mathbin{\sequenceindex{\times}}4}}
\symrow{\begin{gathered}\notationfraction{a}{b}=a\sequenceindex{/}b\\{}=a\mathbin{\sequenceindex{\div}} b\end{gathered}}
{Division: \\ \meaningexplanation{fraction, slash, or division sign ($b\neq0$)}}
{\exmath{\notationfraction{12}{3}&=12\sequenceindex{/}3=12\mathbin{\sequenceindex{\div}}3=4\\\notationfraction{x}{5}&=x\sequenceindex{/}5=x\mathbin{\sequenceindex{\div}}5\\& \qquad \quad =\notationfraction{1}{5}x \> \> \! =0.2x}}
}
\end{minipage}\hfill
\begin{minipage}[t]{.485\linewidth}
\vspace{0pt}
\notationtable{
\symrow{\mathbin{\sequenceindex{\pm}}}{Plus or minus: \meaningexplanation{two choices}}{\exmath{3\mathbin{\sequenceindex{\pm}}1&=4\text{ or }2\\-2\mathbin{\sequenceindex{\pm}}3&=1\text{ or }-5}}
\symrow{\sequenceindex{|}a\sequenceindex{|}}{Absolute value: \meaningexplanation{Distance to $0$}}{\exmath{\sequenceindex{|}-5\sequenceindex{|}&=5\\\sequenceindex{|}3\sequenceindex{|}&=3\quad \sequenceindex{|}0\sequenceindex{|}=0}}
\pairedrowwithcols{l}{a^{\sequenceindex{2}}}{a^{\sequenceindex{3}}}
  {Squared \par \meaningexplanation{Two factors of $a$}}
  {Cubed \par \meaningexplanation{Three factors of $a$}}
  {\begin{array}[t]{@{}r@{}l@{}}3^{\sequenceindex{2}}&{}=3\cdot3=9\\(-2)^{\sequenceindex{2}}&{}=(-2)(-2)=4\end{array}}
  {\begin{array}[t]{@{}r@{}l@{}}2^{\sequenceindex{3}}&{}=2\cdot2\cdot2=8\\(-2)^{\sequenceindex{3}}&{}=(-2)(-2)(-2)\\&{}=-8\end{array}}
\pairedrowwithcols{r@{}E}{\sequenceindex{\sqrt{a}}=a^{\sequenceindex{1/2}}}{\sqrt[\sequenceindex{3}]{a}=a^{\sequenceindex{1/3}}}
  {Square root \quad \> \meaningexplanation{$a\geq0$}}
  {Cube root \qquad \! \meaningexplanation{$a\in\mathbb R$}}
  {\sequenceindex{\sqrt9}&=3\text{ (not }\pm3\text{)}}
  {\sqrt[\sequenceindex{3}]{-8}&=-2}
\pairedrowwithcols{l}{a^{\sequenceindex{n}}}{\sqrt[\sequenceindex{n}]{a}=a^{1/\sequenceindex{n}}}
  {$n$th power}
  {$n$th root}
  {\begin{array}[t]{@{}r@{}l@{}}2^{\sequenceindex{4}}&{}=2\cdot2\cdot2\cdot2=16\\(-2)^{\sequenceindex{5}}&{}=-32\end{array}}
  {\begin{array}[t]{@{}r@{}l@{}}\sqrt[\sequenceindex{4}]{16}&{}=\> \> \> 16^{\frac1{\sequenceindex{4}}} \quad  =2\\\sqrt[\sequenceindex{5}]{-32}&{}=(-32)^{\frac1{\sequenceindex{5}}}=-2\end{array}}
}
\end{minipage}\par\vspace{0mm}
\endgroup
\vspace{3mm}
\compacttables
\renewcommand{\MeaningColumnFraction}{.35}
% Shared heading row keeps both titles and their column rules aligned.
\noindent\begin{minipage}[t]{.49\linewidth}\vspace{0pt}
\sectionheading{\strut Sets}
\end{minipage}\hfill
\begin{minipage}[t]{.49\linewidth}\vspace{0pt}
\sectionheading{\strut Number systems}
\end{minipage}\par\nobreak
\noindent\begin{minipage}[t]{.49\linewidth}\vspace{0pt}
\renewcommand{\ContextHeading}{Example guide}
\examplecontexts{$U$: Universal set \quad $U=\{1,2,3,4\}$ \> | \> $A=\{1,2\}$ \> | \> $B=\{2,3\}$}
% Fit the complement and its note using local cell clearance only.
\setlength{\RowPadding}{1mm}
\notationtable{
\symrow{\mathopen{\sequenceindex{\{}}a,b,c\mathclose{\sequenceindex{\}}}}{\conceptmeaning{Set}{\meaningexplanation{Ignores order and repetition}}}{\exmath{A&=\mathopen{\sequenceindex{\{}}2,1\mathclose{\sequenceindex{\}}}=\mathopen{\sequenceindex{\{}}1,2\mathclose{\sequenceindex{\}}}\\B&=\mathopen{\sequenceindex{\{}}2,3\mathclose{\sequenceindex{\}}}=\mathopen{\sequenceindex{\{}}2,2,3,3\mathclose{\sequenceindex{\}}}}}
\pairedrow{\mathrel{\sequenceindex{\mathord{\in}}}}{\mathrel{\sequenceindex{\mathord{\notin}}}}{Element of}{Not an element of}{2&\mathrel{\sequenceindex{\mathord{\in}}} A \qquad 2 \mathrel{\sequenceindex{\mathord{\in}}} \{2\}}{3&\mathrel{\sequenceindex{\mathord{\notin}}} A \qquad 3 \mathrel{\sequenceindex{\mathord{\notin}}} \{2\}}
\symrow{\{x\mathrel{\sequenceindex{\mathord{:}}}p(x)\}}{\conceptmeaning{Set-builder notation}{\meaningexplanation{Select elements satisfying $p(x)$}}}{\exmath[l]{A=\{x\in U\mathrel{\sequenceindex{\mathord{:}}}x\leq2\}\\B=\{x\in U\mathrel{\sequenceindex{\mathord{:}}}1<x<4\}\\\varnothing=\{x\in U\mathrel{\sequenceindex{\mathord{:}}}x>4\}}}
\pairedrow{\sequenceindex{\varnothing}}{\sequenceindex{|}A\sequenceindex{|}}{\conceptmeaning{Empty set}{\meaningexplanation{No elements}}}{\conceptmeaning{Cardinality}{\meaningexplanation{Number of elements}}}{|\sequenceindex{\varnothing}|&=0}{\sequenceindex{|}A\sequenceindex{|}&=2 \qquad \sequenceindex{|}U\sequenceindex{|}=4}
\pairedrowwithcols{l}{A\mathrel{\sequenceindex{\mathord{\subseteq}}} U}{A\mathrel{\sequenceindex{\mathord{\subset}}} U}{\conceptmeaning{Subset}{\meaningexplanation{All elements included; equality allowed}}}{\conceptmeaning{Proper subset}{\meaningexplanation{Included but not equal}}}{\renewcommand{\arraystretch}{0.8}\begin{array}[t]{@{}l@{}}A\mathrel{\sequenceindex{\mathord{\subseteq}}} A \qquad \>\>\>\> A\mathrel{\sequenceindex{\mathord{\subseteq}}} U\\ \\ \{1\}\mathrel{\sequenceindex{\mathord{\subseteq}}}\{1,2\}\end{array}}{\renewcommand{\arraystretch}{0.8}\begin{array}[t]{@{}l@{}}\{1\}\mathrel{\sequenceindex{\mathord{\subset}}} A \qquad A\mathrel{\sequenceindex{\mathord{\subset}}} U\\ \\ A\mathrel{\sequenceindex{\mathord{\subset}}}\{1,2,3\}\end{array}}
\pairedrow{A\mathbin{\sequenceindex{\cup}} B}{A\mathbin{\sequenceindex{\cap}} B}{\conceptmeaning{Union}{\meaningexplanation{In either or both}}}{\conceptmeaning{Intersection}{\meaningexplanation{In both}}}{A\mathbin{\sequenceindex{\cup}} B&=\{1,2,3\}}{A\mathbin{\sequenceindex{\cap}} B&=\{2\}}
\symrow{A\mathbin{\sequenceindex{\setminus}} B}{\conceptmeaning{Difference}{\meaningexplanation{In $A$, outside $B$}}}{\quad\>\>\! \exmath{A & \mathbin{\sequenceindex{\setminus}} B=\{1\}} \\ \exmath{\{3,4\} &\mathbin{\sequenceindex{\setminus}} \{4\}=\{3\}}}
\symrow{A^{\sequenceindex{c}} \text{ or } A^{\sequenceindex{\prime}}}{\conceptmeaning{Complement}{\meaningexplanation{In $U$, outside $A$}}}{\exmath{A^{\sequenceindex{c}}&=\{3,4\} \quad\>\>\>\>\>\> \{1\}^{\sequenceindex{c}} =  \{2,3,4\}\\B^{\sequenceindex{c}}&=\{1,4\} \quad\>\>\>\>\>\> \{1,2,3\}^{\sequenceindex{c}} =  \{4\}}}
}\par\vspace{1mm}
{\FooterFont \qquad\qquad $2$ is an element of $A$ \qquad | \qquad $\{2\}$ is a subset of $A$\par}
\end{minipage}\hfill
\begin{minipage}[t]{.49\linewidth}\vspace{0pt}
% Transfer three percent of this table from to Meaning.
\renewcommand{\SymbolColumnFraction}{.19}
\renewcommand{\MeaningColumnFraction}{.38}
% Fit the expanded number-system reference without changing typography.
\setlength{\RowPadding}{1mm}
\notationtable{
\pairedrow{\sequenceindex{\mathbb N}}{\sequenceindex{\mathbb Z}}{Naturals}{Integers}{\sequenceindex{\mathbb N}&=\{1,2,3,\ldots\}}{\sequenceindex{\mathbb Z}&=\{\ldots, -2, -1, 0, 1,2,\ldots\}}
\pairedrow{\sequenceindex{\mathbb Q}}{\sequenceindex{\mathbb R}}{\conceptmeaning{Rationals}{\meaningexplanation{$p/q$; $p,q\in\mathbb Z$, $q\ne0$}}}{\conceptmeaning{Reals}{\meaningexplanation{The number line}}}{\tfrac12,-3&\in\sequenceindex{\mathbb Q}}{\sqrt2,\pi,\tfrac12,-3&\in\sequenceindex{\mathbb R}}
\symrow{\mathbb R\setminus\mathbb Q}{Irrationals}{\exmath{\sqrt2,\pi&\in\mathbb R\setminus\mathbb Q}}
\symrow{\sequenceindex{\mathbb C}}{\conceptmeaning{Complex numbers}{\meaningexplanation{$a+bi$; $a,b\in\mathbb R$, $i^2=-1$}}}{\exmath{2+3i&\in\sequenceindex{\mathbb C}}}
\symrow{\begin{gathered}[t]z=\\\sequenceindex{r}(\cos\sequenceindex{\theta}\\{}+i\sin\sequenceindex{\theta})\end{gathered}}
  {\conceptmeaning{Polar form:}{\meaningexplanation{$z=a+bi$:\\ $\sequenceindex{r}=|z|=\sqrt{a^2+b^2}$;\\ $\sequenceindex{\theta}$: argument in radians ($z\ne0$).}}}
  {\exmath[l]{1+i=\sequenceindex{\sqrt{2}}(\cos(\sequenceindex{\pi/4})\\{}+i\sin(\sequenceindex{\pi/4}));\\\sequenceindex{r}=\sequenceindex{\sqrt2},\quad\sequenceindex{\theta}=\sequenceindex{\pi/4}.}}
\pairedrow{\mathbb R^{\sequenceindex{2}}}{\mathbb R^{\sequenceindex{3}}}
  {\conceptmeaning{Real coordinate plane}{\meaningexplanation{Ordered pairs}}}
  {\conceptmeaning{Real coordinate space}{\meaningexplanation{Ordered triples}}}
  {(2,-3)&\in\mathbb R^{\sequenceindex{2}}}{(1,0,-2)&\in\mathbb R^{\sequenceindex{3}}}
\symrow{\mathbb R^{\sequenceindex{+}}}{\conceptmeaning{Positive real numbers}{\meaningexplanation{Excludes zero}}}{\exmath{\mathbb R^{\sequenceindex{+}}&=\{x\in\mathbb R:x>0\}\\&=(0,\infty)}}
\pairedrowwithcols{l@{}E}{\mathopen{\sequenceindex{(}}a,b\mathclose{\sequenceindex{)}}}{\mathopen{\sequenceindex{[}}a,b\mathclose{\sequenceindex{]}}}{\conceptmeaning{Open interval}{\meaningexplanation{Excludes endpoints}}}{\conceptmeaning{Closed interval}{\meaningexplanation{Includes endpoints}}}{x\in\mathopen{\sequenceindex{(}}1,3\mathclose{\sequenceindex{)}}:&\quad 1<x<3}{x\in\mathopen{\sequenceindex{[}}1,3\mathclose{\sequenceindex{]}}:&\quad 1\leq x\leq3}
\pairedrowwithcols{l}{\mathopen{\sequenceindex{[}}a,b\mathclose{\sequenceindex{)}}}{\mathopen{\sequenceindex{(}}a,b\mathclose{\sequenceindex{]}}}{\conceptmeaning{Half-open interval}{\meaningexplanation{Includes left endpoint}}}{\conceptmeaning{Half-open interval}{\meaningexplanation{Includes right endpoint}}}{\begin{array}[t]{@{}l@{}}2\in\mathopen{\sequenceindex{[}}2,5\mathclose{\sequenceindex{)}}\\5\notin\mathopen{\sequenceindex{[}}2,5\mathclose{\sequenceindex{)}}\end{array}}{\begin{array}[t]{@{}l@{}}2\notin\mathopen{\sequenceindex{(}}2,5\mathclose{\sequenceindex{]}}\\5\in\mathopen{\sequenceindex{(}}2,5\mathclose{\sequenceindex{]}}\end{array}}
\pairedrowwithcols{l}{\sequenceindex{\infty}}{-\sequenceindex{\infty}}{Positive infinity}{Negative infinity}{[0,\sequenceindex{\infty})}{(-\sequenceindex{\infty},0]}
}\vspace{1mm}
{\FooterFont Intervals assume $a<b$. Some texts include $0$ in $\mathbb N$\par}
\end{minipage}\par
\par\vspace{3mm}
\compacttables
\sectionheading{Greek symbols and constants}
% One alphabetic letter table beside a separate constants table.
\begingroup
\setlength{\RowPadding}{0.15mm}
\setlength{\tabcolsep}{0.7mm}
\renewcommand{\arraystretch}{0.8}
\arrayrulecolor{rulecolor}
\newcommand{\greekentry}[3]{%
  \paddedcell{\DefinitionFont #1}&%
  \paddedcell{\centering\SymbolFont\ensuremath{#2}}&%
  \paddedcell{\centering\SymbolFont\ensuremath{#3}}%
}
\newcommand{\greekhead}{%
  \paddedcell{\TableHeadFont Name}&%
  \paddedcell{\centering\TableHeadFont Lowercase Symbol}&%
  \paddedcell{\centering\TableHeadFont Uppercase Symbol}%
}
\noindent\begin{minipage}[t]{.49\linewidth}
\vspace{0pt}
\setlength{\tabcolsep}{0.25mm}
\setlength{\RowPadding}{0.9mm}
% Read down the left group, then the right group: six letters each.
\begin{tabularx}{\linewidth}{@{}>{\raggedright\arraybackslash}X>{\centering\arraybackslash}p{.176\linewidth}>{\centering\arraybackslash}p{.176\linewidth}@{\hspace{0.7mm}{\color{lightgray}\vrule width 1.5pt}\hspace{0.7mm}}>{\raggedright\arraybackslash}X>{\centering\arraybackslash}p{.176\linewidth}>{\centering\arraybackslash}p{.176\linewidth}@{}}
\greekhead&\greekhead\\\hline
\greekentry{Alpha}{\alpha}{\mathrm A}&\greekentry{Mu}{\mu}{\mathrm M}\\\hline
\greekentry{Beta}{\beta}{\mathrm B}&\greekentry{Pi}{\pi}{\Pi}\\\hline
\greekentry{Gamma}{\gamma}{\Gamma}&\greekentry{Rho}{\rho}{\mathrm P}\\\hline
\greekentry{Delta}{\delta}{\Delta}&\greekentry{Sigma}{\sigma}{\Sigma}\\\hline
\greekentry{Theta}{\theta}{\Theta}&\greekentry{Phi}{\phi}{\Phi}\\\hline
\greekentry{Lambda}{\lambda}{\Lambda}&\greekentry{Omega}{\omega}{\Omega}\\\hline
\end{tabularx}
\end{minipage}\hfill
\begin{minipage}[t]{.49\linewidth}
\vspace{0pt}
\setlength{\RowPadding}{1.1mm}
\begin{tabularx}{\linewidth}{@{}>{\centering\arraybackslash}p{.16\linewidth}>{\raggedright\arraybackslash}p{.26\linewidth}>{\raggedright\arraybackslash}X@{}}
\paddedcell{\centering\TableHeadFont Constant}&
\paddedcell{\TableHeadFont Name}&
\paddedcell{\TableHeadFont Meaning / Value}\\\hline
\paddedcell{\centering\SymbolFont $\sequenceindex{\pi}$}&
\paddedcell{\DefinitionFont Pi}&
\paddedcell{\DefinitionFont Circle constant\par $\sequenceindex{\pi}=C/d\approx\sequenceindex{3.14159}$\par \meaningexplanation{$C$: circumference; $d$: diameter.}}\\\hline
\paddedcell{\centering\SymbolFont $\sequenceindex{e}$}&
\paddedcell{\DefinitionFont Euler's number}&
\paddedcell{\DefinitionFont Base of natural logarithms\par $\sequenceindex{e}\approx\sequenceindex{2.71828}$\par \meaningexplanation{$\ln x=y\iff \sequenceindex{e}^{y}=x$ ($x>0$);\\ $\ln1=0$, $\ln\sequenceindex{e}=1$.}}\\\hline
\paddedcell{\centering\SymbolFont $\sequenceindex{i}$}&
\paddedcell{\DefinitionFont Imaginary unit}&
\paddedcell{\DefinitionFont $\sequenceindex{i}^2=-1$\par \meaningexplanation{$a+b\sequenceindex{i}$; $a,b\in\mathbb R$.}}\\\hline
\end{tabularx}\par\vspace{0.8mm}
\end{minipage}
\endgroup
{\color{navy}\hrule height 0.5pt}\vspace{0.7mm}
\endgroup
\newpage
\posterheader{A-2\quad FUNCTIONS AND PROBABILITY}
\begingroup
\compacttables
\begingroup
% Compact spacing and balanced columns apply only to this middle poster.
\setlength{\RowPadding}{0.8mm}
\renewcommand{\SymbolColumnFraction}{.24}
\renewcommand{\MeaningColumnFraction}{.38}
\sectionheading{Function notation}
\begingroup
\setlength{\tabcolsep}{0.5mm}
\setlength{\RowPadding}{0.5mm}
\renewcommand{\arraystretch}{1}
\renewcommand{\ExampleLineStretch}{1.02}
\arrayrulecolor{rulecolor}
% Separate notation and definitions beside the three original example columns.
\newcommand{\functionrow}[6]{%
  \paddedcell{\centering\SymbolFont\ensuremath{\displaystyle #1}}&%
  \paddedcell{\DefinitionFont\conceptmeaning{#2}{\meaningexplanation{#3}}}&%
  \paddedcell{\ExampleFont #4}&\paddedcell{\ExampleFont #5}&\paddedcell{\ExampleFont #6}\\\hline}
\begin{tabularx}{\linewidth}{@{}p{.14\linewidth}p{.22\linewidth}p{.20\linewidth}p{.20\linewidth}X@{}}
\paddedcell{\centering\TableHeadFont Symbol}&
\paddedcell{\TableHeadFont Meaning}&
\paddedcell{\TableHeadFont Square function}&
\paddedcell{\TableHeadFont Square-root function}&
\paddedcell{\TableHeadFont Linear function}\\\hline
\multicolumn{2}{@{}p{\dimexpr.36\linewidth+2\tabcolsep\relax}}{%
\paddedcell{\ExampleFont $h$ names the function \par $x$ is an input; $h(x)$ is its output.\par
$(x,h(x))$ is a point on its graph.\par
A function assigns exactly one output in $K$ to each input in $D$.}}&
% Native TikZ coordinates retain the numerical values from Function value below.
% Labels remain 9 pt at the physical 5/6 page size; only local graph units grow.
% Quiet projections are drawn first so curves and marked points stay prominent.
\paddedcell{\ExampleFont $f(x)=x^2$\par
  \begin{tikzpicture}[font=\fontsize{10.8}{12}\selectfont,
    every node/.style={inner sep=0pt},line cap=round,
    guide/.style={draw=navy!45,line width=.25pt,dash pattern=on 1.2pt off 1.2pt},
    leader/.style={draw=navy!65,line width=.3pt}]
    \path[use as bounding box] (0,0) rectangle (56mm,44mm);
    \begin{scope}[shift={(29mm,7mm)},x=5mm,y=2.2mm]
      \draw[guide] (-2,0)--(-2,4)--(0,4);
      \draw[guide] (3,0)--(3,9)--(0,9);
      \draw[->,line width=.35pt] (-3.4,0)--(4.6,0);
      \draw[->,line width=.35pt] (0,-.5)--(0,11);
      \draw[coral,line width=.7pt,domain=-3.1:3.4,samples=65]
        plot (\x,{\x*\x});
      \fill (-2,4) circle[radius=1.1pt] (3,9) circle[radius=1.1pt];
      \node[below=1.5mm] at (-2,0) {$-2$};
      \node[below=1.5mm] at (3,0) {$3$};
      \node[right=1.5mm] at (0,4) {$4$};
      \node[left=1.5mm] at (0,9) {$9$};
      \coordinate (fNegativePoint) at (-2,4);
      \coordinate (fPositivePoint) at (3,9);
    \end{scope}
    \node[anchor=center] (fNegativeLabel) at (19mm,37mm) {$(-2,f(-2))$};
    \node[anchor=center] (fPositiveLabel) at (37.5mm,33mm) {$(3,f(3))$};
    \draw[leader,shorten <=1.5pt,shorten >=1mm]
      (fNegativePoint)--(fNegativeLabel);
    \draw[leader,shorten <=1.5pt,shorten >=1mm]
      (fPositivePoint)--(fPositiveLabel);
  \end{tikzpicture}}&
\paddedcell{\ExampleFont $g(x)=\sqrt{x}$\par
  \begin{tikzpicture}[font=\fontsize{10.8}{12}\selectfont,
    every node/.style={inner sep=0pt},line cap=round,
    guide/.style={draw=navy!45,line width=.25pt,dash pattern=on 1.2pt off 1.2pt},
    leader/.style={draw=navy!65,line width=.3pt}]
    \path[use as bounding box] (0,0) rectangle (56mm,44mm);
    \begin{scope}[shift={(6mm,7mm)},x=2.5mm,y=5.5mm]
      \draw[guide] (9,0)--(9,3)--(0,3);
      \draw[guide] (16,0)--(16,4)--(0,4);
      \draw[->,line width=.35pt] (-.5,0)--(19,0);
      \draw[->,line width=.35pt] (0,-.2)--(0,5.2);
      % Parameterization x=t^2, y=t gives a smooth start, only for x>=0.
      \draw[coral,line width=.7pt,domain=0:4.25,samples=65,variable=\t]
        plot ({\t*\t},\t);
      \fill (9,3) circle[radius=1.1pt] (16,4) circle[radius=1.1pt];
      \node[below=1.5mm] at (9,0) {$9$};
      \node[below=1.5mm] at (16,0) {$16$};
      \node[left=1.5mm] at (0,3) {$3$};
      \node[left=1.5mm] at (0,4) {$4$};
      \coordinate (gNinePoint) at (9,3);
      \coordinate (gSixteenPoint) at (16,4);
    \end{scope}
    \node[anchor=center] (gNineLabel) at (37.25mm,17mm)
      {$(9,g(9))$};
    \node[anchor=center] (gSixteenLabel) at (35mm,35mm)
      {$(16,g(16))$};
    \draw[leader,shorten <=1.5pt,shorten >=1mm]
      (gNinePoint)--(gNineLabel);
    \draw[leader,shorten <=1.5pt,shorten >=1mm]
      (gSixteenPoint)--(gSixteenLabel);
  \end{tikzpicture}}&
\paddedcell{\ExampleFont $\beta(x)=2x+1$\par
  \begin{tikzpicture}[font=\fontsize{10.8}{12}\selectfont,
    every node/.style={inner sep=0pt},line cap=round,
    guide/.style={draw=navy!45,line width=.25pt,dash pattern=on 1.2pt off 1.2pt},
    leader/.style={draw=navy!65,line width=.3pt}]
    \path[use as bounding box] (0,0) rectangle (56mm,44mm);
    \begin{scope}[shift={(25mm,16mm)},x=6mm,y=2.7mm]
      \draw[guide] (-1,0)--(-1,-1)--(0,-1);
      \draw[guide] (3,0)--(3,7)--(0,7);
      \draw[->,line width=.35pt] (-3.6,0)--(4.2,0);
      \draw[->,line width=.35pt] (0,-2.2)--(0,8.5);
      \draw[coral,line width=.7pt,domain=-1.6:3.7]
        plot (\x,{2*\x+1});
      \fill (-1,-1) circle[radius=1.1pt] (3,7) circle[radius=1.1pt];
      \node[above=1.5mm] at (-1,0) {$-1$};
      \node[below=1.5mm] at (3,0) {$3$};
      \node[right=1.5mm] at (0,-1) {$-1$};
      \node[left=1.5mm] at (0,7) {$7$};
      \coordinate (betaNegativePoint) at (-1,-1);
      \coordinate (betaPositivePoint) at (3,7);
    \end{scope}
    % Convert the former corner coordinates to centers using the label dimensions.
    \setbox0=\hbox{\fontsize{10.8}{12}\selectfont $(-1,\beta(-1))$}%
    \node[anchor=center] (betaNegativeLabel) at
      (\dimexpr23mm-.5\wd0-.5\pgflinewidth\relax,
       \dimexpr29mm-.5\ht0-.5\dp0-.5\pgflinewidth\relax)
      {$(-1,\beta(-1))$};
    \setbox0=\hbox{\fontsize{10.8}{12}\selectfont $(3,\beta(3))$}%
    \node[anchor=center] (betaPositiveLabel) at
      (\dimexpr28mm+.5\wd0+.5\pgflinewidth\relax,
       \dimexpr38mm+.5\ht0+.5\dp0+.5\pgflinewidth\relax)
      {$(3,\beta(3))$};
    \draw[leader,shorten <=1mm,shorten >=1.5pt]
      (betaNegativeLabel.south)--(betaNegativeLabel.south |- 0,11mm)--(betaNegativePoint);
    \draw[leader,shorten <=1.5pt,shorten >=1mm]
      (betaPositivePoint)--(betaPositiveLabel);
  \end{tikzpicture}}\\\hline
\functionrow{h:D\mathrel{\sequenceindex{\mathord{\to}}} K}{Function mapping}{Domain $D$ \quad codomain $K$}
  {$f:\mathbb R\mathrel{\sequenceindex{\mathord{\to}}}\mathbb R$}
  {$g:[0,\infty)\mathrel{\sequenceindex{\mathord{\to}}}\mathbb R$}
  {$\beta:\mathbb R\mathrel{\sequenceindex{\mathord{\to}}}\mathbb R$}
\functionrow{\mathop{\sequenceindex{\mathord{\operatorname{dom}}}}(h)=\sequenceindex{D}}{Domain}{Allowed inputs}
  {$\mathop{\sequenceindex{\mathord{\operatorname{dom}}}}(f)=\sequenceindex{\mathbb R}$}
  {$\mathop{\sequenceindex{\mathord{\operatorname{dom}}}}(g)=\sequenceindex{[0,\infty)}$}
  {$\mathop{\sequenceindex{\mathord{\operatorname{dom}}}}(\beta)=\sequenceindex{\mathbb R}$}
\functionrow{\mathop{\sequenceindex{\mathord{\operatorname{codom}}}}(h)=\sequenceindex{K}}{Codomain}{Declared output target}
  {$\mathop{\sequenceindex{\mathord{\operatorname{codom}}}}(f)=\sequenceindex{\mathbb R}$}
  {$\mathop{\sequenceindex{\mathord{\operatorname{codom}}}}(g)=\sequenceindex{\mathbb R}$}
  {$\mathop{\sequenceindex{\mathord{\operatorname{codom}}}}(\beta)=\sequenceindex{\mathbb R}$}
\functionrow{x\mathrel{\sequenceindex{\mathord{\mapsto}}} h(x)}{Mapping rule}{$\mapsto$ maps inputs\enspace $\to$ links sets}
  {\quad\>\!\! $x\mathrel{\sequenceindex{\mathord{\mapsto}}} x^2$}
  {\quad\>\!\! $x\mathrel{\sequenceindex{\mathord{\mapsto}}}\sqrt{x}$}
  {\quad\>\!\! $x\mathrel{\sequenceindex{\mathord{\mapsto}}} 2x+1$}
\functionrow{\sequenceindex{h(x)}}{Function value}{Output at input $x$}
  {\exmath[l]{f\mathopen{\sequenceindex{(}}3\mathclose{\sequenceindex{)}}=3^2=\sequenceindex{9}\\f\mathopen{\sequenceindex{(}}-2\mathclose{\sequenceindex{)}}=(-2)^2=\sequenceindex{4}}}
  {\exmath[l]{g\mathopen{\sequenceindex{(}}9\mathclose{\sequenceindex{)}}=\sqrt{9}=\sequenceindex{3}\\g\mathopen{\sequenceindex{(}}16\mathclose{\sequenceindex{)}}=\sqrt{16}=\sequenceindex{4}}}
  {\exmath[l]{\beta\mathopen{\sequenceindex{(}}3\mathclose{\sequenceindex{)}}=2(3)+1=\sequenceindex{7}\\\beta\mathopen{\sequenceindex{(}}-1\mathclose{\sequenceindex{)}}=2(-1)+1=\sequenceindex{-1}}}
\functionrow{\mathop{\sequenceindex{\mathord{\operatorname{range}}}}(h)}{Range}{Only attained outputs}
  {$\mathop{\sequenceindex{\mathord{\operatorname{range}}}}(f)=\sequenceindex{[0,\infty)}$}
  {$\mathop{\sequenceindex{\mathord{\operatorname{range}}}}(g)=\sequenceindex{[0,\infty)}$}
  {$\mathop{\sequenceindex{\mathord{\operatorname{range}}}}(\beta)=\sequenceindex{\mathbb R}$}
\end{tabularx}\par
\endgroup
\vspace{2mm}
\sectionheading{Probability and events}
\begingroup
\setlength{\tabcolsep}{0.6mm}
% Share exact column dimensions between the context panel and main table.
\setlength{\SymbolWidth}{.13\linewidth}
\setlength{\MeaningWidth}{.25\linewidth}
\setlength{\ExampleWidth}{\dimexpr(\linewidth-\SymbolWidth-\MeaningWidth-6\tabcolsep)/2\relax}
\newcolumntype{V}{@{}p{\SymbolWidth}p{\MeaningWidth}p{\ExampleWidth}p{\ExampleWidth}@{}}
\begingroup
\setlength{\fboxsep}{0.5mm}
\setlength{\fboxrule}{0.35pt}
\noindent\fcolorbox{coral}{coral!10}{%
  \begin{minipage}{\dimexpr\linewidth-2\fboxsep-2\fboxrule\relax}
  \ExampleFont\color{navy}%
  % Cancel the frame inset to retain the main table's column origins.
  \makebox[\linewidth][l]{%
  \hspace*{\dimexpr-\fboxsep-\fboxrule\relax}%
  \begin{tabular}{V}
  \multicolumn{2}{@{}p{\dimexpr\SymbolWidth+\MeaningWidth+2\tabcolsep\relax}}{%
    \begin{minipage}[t]{\linewidth}\raggedright
    {\centering\textbf{Example contexts}\par}
    \vspace{0.5mm}
    \textbf{Fractions, decimals, and percentages can express the same probability:}\par
    \vspace{0.5mm}
    {\centering
    $\displaystyle P(A)=\frac{1}{2}=0.5=50\%$\par
    \vspace{0.5mm}
    $\displaystyle P(F)=\frac{3}{4}=0.75=75\%$\par}
    \vspace{0.5mm}
    Divide the numerator by the denominator to obtain the decimal representation.
    Multiply the decimal by $100$ to obtain the percentage.\par
    \end{minipage}}&
  \begin{minipage}[t]{\ExampleWidth}\raggedright
  \textbf{Die: one fair die roll}\par
  A six-sided die\par
  $\Omega=\{1,2,3,4,5,6\}$\par
  \begingroup
  \let\tabularxcolumn\centeredtextcolumn
  \begin{tabularx}{\linewidth}{@{}>{\collectcell\contextcell}l<{\endcollectcell}>{\collectcell\subrowcell}X<{\endcollectcell}@{}}
  \textbf{Symbol / event set}&\textbf{Meaning}\\
  $A=\{1,3,5\}$&Odd result\\
  $B=\{1,2\}$&Result at most $2$\\
  $C=\{2,4,6\}$&Even result
  \end{tabularx}\endgroup
  \end{minipage}&
  \begin{minipage}[t]{\ExampleWidth}\raggedright
  \textbf{Coins: two independent fair tosses}\par
  $\mathrm H$: heads, $\mathrm T$: tails\par
  $\Omega=\{\mathrm{HH},\mathrm{HT},\mathrm{TH},\mathrm{TT}\}$\par
  \begingroup
  \let\tabularxcolumn\centeredtextcolumn
  \begin{tabularx}{\linewidth}{@{}>{\collectcell\contextcell}l<{\endcollectcell}>{\collectcell\subrowcell}X<{\endcollectcell}@{}}
  \textbf{Symbol / event set}&\textbf{Meaning}\\
  $F=\{\mathrm{HH},\mathrm{HT},\mathrm{TH}\}$&At least one head\\
  $G=\{\mathrm{HH},\mathrm{TT}\}$&Matching results\\
  $J=\{\mathrm{TT}\}$&No heads
  \end{tabularx}\endgroup
  \end{minipage}
  \end{tabular}}\par
  \end{minipage}}\par
\endgroup
\vspace{0.7mm}
% Compact this table locally so the complete probability poster stays on one page.
\setlength{\RowPadding}{0.15mm}
\renewcommand{\arraystretch}{1}
\renewcommand{\ExampleLineStretch}{1.02}
\arrayrulecolor{rulecolor}
% Separate notation and definitions beside the two original experiment columns.
\newcommand{\probrow}[5]{%
  \paddedcell{\centering\SymbolFont\ensuremath{\displaystyle #1}}&%
  \paddedcell{\DefinitionFont\conceptmeaning{\if\relax\detokenize{#2}\relax\else#2:\fi}{\meaningexplanation{#3}}}&%
  \paddedcell{\ExampleFont #4}&\paddedcell{\ExampleFont #5}\\\hline}
\begin{tabular}{V}
\paddedcell{\centering\TableHeadFont Symbol}&
\paddedcell{\TableHeadFont Meaning}&
\paddedcell{\TableHeadFont Die examples}&
\paddedcell{\TableHeadFont Coins examples}\\\hline
\probrow{\sequenceindex{\Omega}\quad\sequenceindex{\omega}}{Sample space; outcome}{All outcomes; one outcome}{$|\sequenceindex{\Omega}|=6$ \qquad $\sequenceindex{\omega}=3\in\sequenceindex{\Omega}$\par \qquad \qquad \quad \> $\sequenceindex{\omega}=6\in\sequenceindex{\Omega}$}{$|\sequenceindex{\Omega}|=4$ \qquad $\sequenceindex{\omega}=\mathrm{HT}\in\sequenceindex{\Omega}$\par \qquad \qquad \quad \> $\sequenceindex{\omega}=\mathrm{HH}\in\sequenceindex{\Omega}$}
\probrow{\sequenceindex{E}\subseteq\Omega}{Input: event}{Outcome set, not a probability}{\exmath[l]{\sequenceindex{E}_1=\{1,4,6\},\quad |\sequenceindex{E}_1|=3\\\sequenceindex{E}_2=\{2,5\},\quad |\sequenceindex{E}_2|=2}}{\exmath[l]{\sequenceindex{E}_1=\{\mathrm{HH},\mathrm{HT}\},\quad |\sequenceindex{E}_1|=2\\\sequenceindex{E}_2=\{\mathrm{HH},\mathrm{TH}\},\quad |\sequenceindex{E}_2|=2}}
\probrow{\begin{gathered}\sequenceindex{\mathcal P}(\Omega)\\{}=\{E:E\subseteq\Omega\}\end{gathered}}{Power set; domain of $P$}{All events (subsets of $\Omega$).}{$|\sequenceindex{\mathcal P}(\Omega)|=2^6=64$; e.g.\par $\varnothing,\Omega,A,B,C,\{1\},A\cap B$}{$|\sequenceindex{\mathcal P}(\Omega)|=2^4=16$; e.g.\par $\varnothing,\Omega,F,G,J,\{\mathrm{HT}\},F\cap G$}
\probrow{\sequenceindex{P}}{Probability measure}{Function $\sequenceindex{P}:\mathcal P(\Omega)\to[0,1]$, $E\mapsto \sequenceindex{P}(E)$}{\exmath[l@{\quad}l]{\sequenceindex{P}(\varnothing)=0&\sequenceindex{P}(A)=\tfrac12\\\sequenceindex{P}(B)=\tfrac13&\sequenceindex{P}(\{6\})=\tfrac16}}{\exmath[l@{\quad}l]{\sequenceindex{P}(\Omega)=1&\sequenceindex{P}(F)=\tfrac34\\\sequenceindex{P}(G)=\tfrac12&\sequenceindex{P}(\{\mathrm{HT}\})=\tfrac14}}
\probrow{[0,1]}{Codomain}{Declared numerical target }{$[0,1]$ or $[0\%,100\%]$}{$[0,1]$ $[0,1]$ or $[0\%,100\%]$}

\probrow{\sequenceindex{P(E)}}{Output: probability}{Numerical value; $P(\varnothing)=\sequenceindex{0}$}{$P(A)=\sequenceindex{\tfrac12}$ \quad $P(B)=\sequenceindex{\tfrac13}$ \quad $P(C)=\sequenceindex{\tfrac12}$}{$P(F)=\sequenceindex{\tfrac34}$ \quad $P(G)=\sequenceindex{\tfrac12}$ \quad $P(J)=\sequenceindex{\tfrac14}$}
\probrow{\sequenceindex{\text{Range}}(P)}{}{Attained probabilities}{$\sequenceindex{\{0,\tfrac16,\tfrac26,\ldots,1\}}$}{$\sequenceindex{\{0,\tfrac14,\tfrac12,\tfrac34,1\}}$}
\probrow{P(E)=\frac{\sequenceindex{|}E\sequenceindex{|}}{\sequenceindex{|}\Omega\sequenceindex{|}}}{Classical probability}{Equally likely finite outcomes:}{\exmath[l]{A=\{1,3,5\}\\P(A)=\frac{\sequenceindex{|}A\sequenceindex{|}}{\sequenceindex{|}\Omega\sequenceindex{|}}=\frac{\sequenceindex{3}}{\sequenceindex{6}}=\tfrac12}}{\exmath[l]{F=\{\mathrm{HH},\mathrm{HT},\mathrm{TH}\}\\P(F)=\frac{\sequenceindex{|}F\sequenceindex{|}}{\sequenceindex{|}\Omega\sequenceindex{|}}=\frac{\sequenceindex{3}}{\sequenceindex{4}}}}
\probrow{E^{\sequenceindex{c}}}{Complement}{In $\Omega$, outside $E$; $P(E^{\sequenceindex{c}})=1-P(E)$}{\exmath[l@{\quad}l]{A^{\sequenceindex{c}}=C&\{1,2\}^{\sequenceindex{c}}=\{3,4,5,6\}\\P(A^{\sequenceindex{c}})=\tfrac12&P(\{1,2\}^{\sequenceindex{c}})=\tfrac46=\tfrac23}}{\exmath[l@{\quad}l]{F^{\sequenceindex{c}}=J&\{\mathrm{HT}\}^{\sequenceindex{c}}=\{\mathrm{HH},\mathrm{TH},\mathrm{TT}\}\\P(F^{\sequenceindex{c}})=\tfrac14&P(\{\mathrm{HT}\}^{\sequenceindex{c}})=\tfrac34}}
\probrow{E_1\mathbin{\sequenceindex{\cup}} E_2}{Union}{$E_1$ or $E_2$, or both}{$A\mathbin{\sequenceindex{\cup}} B=A+B=\{1,2,3,5\}$}{$F\mathbin{\sequenceindex{\cup}} G=F+G=\Omega$}
\probrow{E_1\mathbin{\sequenceindex{\cap}} E_2}{Intersection}{$E_1$ and $E_2$ together}{$A\mathbin{\sequenceindex{\cap}} B=AB=\{1\}$}{$F\mathbin{\sequenceindex{\cap}} G=FG=\{\mathrm{HH}\}$}
\probrow{\sequenceindex{E_1\cap E_2=\varnothing}}{Disjoint}{Mutually exclusive:}{$\sequenceindex{A\cap C= AC = \varnothing}$}{$\sequenceindex{F\cap J=\varnothing}$}
\probrow{}{Addition for disjoint events}{If $E_1\cap E_2=\varnothing$ \qquad then $P(E_1\cup E_2)=P(E_1)+P(E_2)$}{$P(A\cup C)=\tfrac12+\tfrac12=1$}{$P(F\cup J)=\tfrac34+\tfrac14=1$}
\paddedcell{\centering\SymbolFont $\displaystyle
  \begin{gathered}P(E_1\cap E_2)\\{}=P(E_1)P(E_2)\end{gathered}$}&
\paddedcell{\DefinitionFont\conceptmeaning{Independent events}{%
  \meaningexplanation{The product rule at left defines independence.}\par
  \meaningexplanation{Knowing one does not change the other's probability.}}}&
\multicolumn{2}{p{\dimexpr2\ExampleWidth+2\tabcolsep\relax}@{}}{\paddedcell{\ExampleFont\raggedright
  \textbf{Coin example: independence}\par
  \begingroup
  \let\tabularxcolumn\centeredtextcolumn
  \begin{tabularx}{\linewidth}{@{}>{\collectcell\contextcell}l<{\endcollectcell}>{\collectcell\subrowcell}X<{\endcollectcell}@{}}
  \textbf{Symbol}&\textbf{Meaning}\\
  $E_1=\{\mathrm{HH},\mathrm{HT}\}$&First toss is heads\\
  $E_2=\{\mathrm{HH},\mathrm{TH}\}$&Second toss is heads\\
  $E_1\cap E_2=\{\mathrm{HH}\}$&Both tosses are heads\\
  $P(E_1)=P(E_2)=\tfrac12$&Each event has two of four outcomes\\
  $P(E_1\cap E_2)=P(E_1)P(E_2)=\tfrac12\cdot\tfrac12=\tfrac14$ \hspace{3mm} &One of four outcomes
  \end{tabularx}\endgroup
}}\\\hline
\end{tabular}\par
\endgroup
{\color{navy}\hrule height 0.5pt}\vspace{0.7mm}
\endgroup
\endgroup
\newpage
\PosterGroup{B}
\posterheader{B-1\quad SEQUENCES AND SUMMATION}
\begingroup
\compacttables
\sectionheading{Sequences and subscripts}
\renewcommand{\ContextHeading}{Reading guide}
\examplecontexts{A \textbf{sequence} is an ordered list of indexed terms. Parentheses show sequences; braces show sets. Here all indices start at $i=1$, and $n\in\mathbb N$, $n\geq1$. Read each row across to compare the three sequences.}
\begingroup
\setlength{\tabcolsep}{0.8mm}
\setlength{\RowPadding}{0.95mm}
\renewcommand{\arraystretch}{1}
\arrayrulecolor{rulecolor}
% Local row helpers serve sequences and partial sums.
% Concept-only rows leave Symbol empty rather than introduce new notation.
% Coral marks isolate newly introduced symbols, indices, and bounds.
% Carry each row's role through every example, including substituted indices and values.
\newcommand{\sequencerow}[6]{%
  \paddedcell{\centering\SymbolFont\ensuremath{\displaystyle #1}}&%
  \paddedcell{\DefinitionFont\conceptmeaning{#2}{\meaningexplanation{#3}}}&%
  \paddedcell{\ExampleFont #4}&\paddedcell{\ExampleFont #5}&\paddedcell{\ExampleFont #6}\\\hline}
\newcommand{\sequencehead}{%
  \paddedcell{\centering\TableHeadFont Symbol}&%
  \paddedcell{\TableHeadFont Meaning}&%
  \paddedcell{\TableHeadFont Sequence $a$}&%
  \paddedcell{\TableHeadFont Sequence $b$}&%
  \paddedcell{\TableHeadFont Sequence $c$}\\\hline}
\begin{tabularx}{\linewidth}{@{}p{.10\linewidth}p{.18\linewidth}p{.225\linewidth}p{.225\linewidth}X@{}}
\sequencehead
\sequencerow{k_{\sequenceindex{i}}=f(\sequenceindex{i})}{Rule and first terms}{Substitute the index.}
  {\exmath[l]{a_{\sequenceindex{i}}=2\cdot\sequenceindex{i}\\(2\cdot\sequenceindex{1},2\cdot\sequenceindex{2},2\cdot\sequenceindex{3},2\cdot\sequenceindex{4},\ldots)\\=(2,4,6,8,\ldots)}}
  {\exmath[l]{b_{\sequenceindex{i}}=\sequenceindex{i}^2\\((\sequenceindex{1})^2,(\sequenceindex{2})^2,(\sequenceindex{3})^2,(\sequenceindex{4})^2,\ldots)\\=(1,4,9,16,\ldots)}}
  {\exmath[l]{c_{\sequenceindex{i}}=2^{\sequenceindex{i}-1}\\(2^{\sequenceindex{1}-1},2^{\sequenceindex{2}-1},2^{\sequenceindex{3}-1},2^{\sequenceindex{4}-1},\ldots)\\=(1,2,4,8,\ldots)}}
\sequencerow{k_{\sequenceindex{i}}}{Indexed term}{Subscript: position.}
  {\exmath[l]{a_{\sequenceindex{3}}=2(\sequenceindex{3})=6\\a_{\sequenceindex{4}}=2(\sequenceindex{4})=8}}
  {\exmath[l]{b_{\sequenceindex{3}}=(\sequenceindex{3})^2=9\\b_{\sequenceindex{5}}=(\sequenceindex{5})^2=25}}
  {\exmath[l]{c_{\sequenceindex{3}}=2^{\sequenceindex{3}-1}=4\\c_{\sequenceindex{6}}=2^{\sequenceindex{6}-1}=32}}
\sequencerow{k_{\sequenceindex{1}},\ k_{\sequenceindex{n}}}{First term; $n$th term}{Term at $1$ or $n$.}
  {$a_{\sequenceindex{1}}=2$ \par $a_{\sequenceindex{n}}=2\sequenceindex{n}$}
  {$b_{\sequenceindex{1}}=1$ \par $b_{\sequenceindex{n}}=\sequenceindex{n}^2$}
  {$c_{\sequenceindex{1}}=1$ \par $c_{\sequenceindex{n}}=2^{\sequenceindex{n}-1}$}
\sequencerow{(k_{i})_{i=1}^{\sequenceindex{4}}}{Finite sequence}{Last index fixes length.}
  {\exmath[l]{(a_{i})_{i=1}^{\sequenceindex{4}}=(2,4,6,8)\\(a_{i})_{i=1}^{\sequenceindex{3}}=(2,4,6)}}
  {\exmath[l]{(b_{i})_{i=1}^{\sequenceindex{4}}=(1,4,9,16)\\(b_{i})_{i=1}^{\sequenceindex{2}}=(1,4)}}
  {\exmath[l]{(c_{i})_{i=1}^{\sequenceindex{4}}=(1,2,4,8)\\(c_{i})_{i=1}^{\sequenceindex{1}}=(1)}}
\sequencerow{(k_{i})_{i=1}^{\sequenceindex{\infty}}}{Infinite sequence, etc.}{No last term; $\infty$ is no index.}
  {$(a_{i})_{i=1}^{\sequenceindex{\infty}}$\par $=(2,4,6,8,\ldots)$}
  {$(b_{i})_{i=1}^{\sequenceindex{\infty}}$\par $=(1,4,9,16,\ldots)$}
  {$(c_{i})_{i=1}^{\sequenceindex{\infty}}$\par $=(1,2,4,8,\ldots)$}
\sequencerow{}{Sequence order}{Order matters.}
  {$(a_{1},a_{2})\ne(a_{2},a_{1})$\par $(2,4)\ne(4,2)$}
  {$(b_{1},b_{2})\ne(b_{2},b_{1})$\par $(1,4)\ne(4,1)$}
  {$(c_{1},c_{2})\ne(c_{2},c_{1})$\par $(1,2)\ne(2,1)$}
\sequencerow{}{Sequence repetition}{Repetition matters.}
  {$(2,2,4)\ne(2,4)$}
  {$(1,1,4)\ne(1,4)$}
  {$(1,1,2)\ne(1,2)$}
\end{tabularx}\par
\vspace{0.5mm}
{\FooterFont An index is not multiplication or an exponent. Sets ignore order and repetition sequences do not.\par}
\vspace{3mm}
\sectionheading{Summation and partial sums}
\renewcommand{\ContextHeading}{From terms to their sum}
\examplecontexts{A \textbf{series} is a sum of sequence terms. The symbol $\mathop{\sequenceindex{\mathord{\sum}}}\displaylimits$ adds the summand as its index runs through consecutive integers between the bounds, \textbf{inclusive}. A partial sum adds finitely many terms.\par
In $S_{n}^{(\sequenceindex{a})}$, $S_{m}^{(\sequenceindex{b})}$, and $S_{l}^{(\sequenceindex{c})}$, the parenthesized superscripts \textbf{label the sequence being summed; they are not powers}. The subscript $\sequenceindex{n}$, $\sequenceindex{m}$, or $\sequenceindex{l}$ counts the terms when the sum starts at $1$.}
\begingroup
\begin{tabularx}{\linewidth}{@{}p{.10\linewidth}p{.2\linewidth}p{.225\linewidth}p{.225\linewidth}X@{}}
\sequencehead
\sequencerow{\sequenceindex{k_{i}} = f(i)}{Summand}{Term being added.}
  {$\sequenceindex{a_{i}}=2i$}
  {$\sequenceindex{b_{i}}=i^2$}
  {$\sequenceindex{c_{i}}=2^{i-1}$}
\sequencerow{\sum_{\sequenceindex{i}=\sequenceindex{1}}^{\sequenceindex{n}}f(\sequenceindex{i})}{Bounds and index}{Lower bound $1$; upper bound $n,m,l$; local index $i$.}
  {$\sum_{\sequenceindex{i}=\sequenceindex{1}}^{\sequenceindex{n}}a_{\sequenceindex{i}}$\par\vspace{1.8mm}Start: $a_{\sequenceindex{1}}$ | End: $a_{\sequenceindex{n}}$.}
  {$\sum_{\sequenceindex{i}=\sequenceindex{1}}^{\sequenceindex{m}}b_{\sequenceindex{i}}$\par\vspace{1.8mm}Start: $b_{\sequenceindex{1}}$ | End: $b_{\sequenceindex{m}}$.}
  {$\sum_{\sequenceindex{i}=\sequenceindex{1}}^{\sequenceindex{l}}c_{\sequenceindex{i}}$\par\vspace{1.8mm}Start: $c_{\sequenceindex{1}}$ | End: $c_{\sequenceindex{l}}$.}
\sequencerow{\sequenceindex{S}_{n}}{Partial sum of the first $n$ terms}{Add through index $n$, $m$, or $l$.}
  {$\sequenceindex{S}_{n}^{(a)}=\sum_{i=1}^{n}a_{i}$\par\vspace{1.8mm}\qquad\! $=a_{1}+\cdots+a_{n}$}
  {$\sequenceindex{S}_{m}^{(b)}=\sum_{i=1}^{m}b_{i}$\par\vspace{1.8mm}\qquad\! $=b_{1}+\cdots+b_{m}$}
  {$\sequenceindex{S}_{l}^{(c)}=\sum_{i=1}^{l}c_{i}$\par\vspace{1.8mm}\qquad\! $=c_{1}+\cdots+c_{l}$}
\sequencerow{\sequenceindex{a}_{4}\text{ vs }\sequenceindex{S}_{4}}{Fourth term versus fourth partial sum}{One term vs. four-term sum.}
  {$\sequenceindex{a}_{4}=\sequenceindex{8}$\par\vspace{1.8mm} $\sequenceindex{S}_{4}^{(a)}=2+4+6+8=\sequenceindex{20}$}
  {$\sequenceindex{b}_{4}=\sequenceindex{16}$\par\vspace{1.8mm} $\sequenceindex{S}_{4}^{(b)}=1+4+9+16=\sequenceindex{30}$}
  {$\sequenceindex{c}_{4}=\sequenceindex{8}$\par\vspace{1.8mm} $\sequenceindex{S}_{4}^{(c)}=1+2+4+8=\sequenceindex{15}$}
\sequencerow{\sequenceindex{i},\sequenceindex{j},\sequenceindex{k}}{Index renaming}{Rename the index throughout; keep bounds.}
  {\exmath[l]{S_{4}^{(a)}=\sum_{\sequenceindex{i}=1}^{4}a_{\sequenceindex{i}}=20\\{}\qquad\!=\sum_{\sequenceindex{j}=1}^{4}a_{\sequenceindex{j}}=20\\{}\qquad\!=\sum_{\sequenceindex{k}=1}^{4}a_{\sequenceindex{k}}=20}}
  {\exmath[l]{S_{4}^{(b)}=\sum_{\sequenceindex{i}=1}^{4}b_{\sequenceindex{i}}=30\\{}\qquad\!=\sum_{\sequenceindex{j}=1}^{4}b_{\sequenceindex{j}}=30\\{}\qquad=\sum_{\sequenceindex{k}=1}^{4}b_{\sequenceindex{k}}=30}}
  {\exmath[l]{S_{4}^{(c)}=\sum_{\sequenceindex{i}=1}^{4}c_{\sequenceindex{i}}=15\\{}\qquad\!=\sum_{\sequenceindex{j}=1}^{4}c_{\sequenceindex{j}}=15\\{}\qquad\!=\sum_{\sequenceindex{k}=1}^{4}c_{\sequenceindex{k}}=15}}
\end{tabularx}\par
\endgroup
\vspace{3mm}
\sectionheading{Worked summation examples}
% Label only the first expansion. Reserve the full stack depth, including
% clearance for highlight padding, a short upward arrow, and a readable label.
\newcommand{\indexedterm}[3]{\mathord{\vtop{\offinterlineskip
  \halign{\hfil##\hfil\cr
    $\mathsurround=0pt #1$\cr
    \noalign{\kern3.5pt}$\mathsurround=0pt\scriptscriptstyle\uparrow$\cr
    \noalign{\kern1.5pt}$\mathsurround=0pt\scriptstyle\sequenceindex{#2=#3}$\cr}}}}
\setlength{\RowPadding}{1.5mm}
\renewcommand{\ExampleLineStretch}{1.3}
\renewcommand{\SymbolColumnFraction}{.23}
\renewcommand{\MeaningColumnFraction}{.35}
\paralleltables{
\symrow{\displaystyle\sum_{\sequenceindex{i=2}}^{4}3}{\conceptmeaning{Constant sum}{\meaningexplanation{Add $3$ for $\sequenceindex{i=2,3,4}$}}}{\exmath[l]{\indexedterm{3}{i}{2}+\indexedterm{3}{i}{3}+\indexedterm{3}{i}{4}=9}}
\symrow{\displaystyle\sum_{\sequenceindex{k=1}}^{3}k^{\sequenceindex{2}}}{\conceptmeaning{Sum of squares}{\meaningexplanation{Square $k$ for\\ $\sequenceindex{k=1,2,3}$}}}{\exmath[l]{\indexedterm{1^{\sequenceindex{2}}}{k}{1}+\indexedterm{2^{\sequenceindex{2}}}{k}{2}+\indexedterm{3^{\sequenceindex{2}}}{k}{3}\\[1.8mm]{}=1+4+9=14}}
\symrow{\displaystyle\sum_{\sequenceindex{j=0}}^{3}2^{j}}{\conceptmeaning{Starting at zero}{\meaningexplanation{Powers of $2$;\\ $\sequenceindex{j=0,1,2,3}$}}}{\exmath[l]{\indexedterm{2^{\sequenceindex{0}}}{j}{0}+\indexedterm{2^{1}}{j}{1}+\indexedterm{2^{2}}{j}{2}+\indexedterm{2^{3}}{j}{3}\\[1.8mm]{}=1+2+4+8=15}}
}{
\symrow{\displaystyle\sum_{\sequenceindex{i=3}}^{5}i}{\conceptmeaning{Shifted lower bound}{\meaningexplanation{Add $i$ for $\sequenceindex{i=3,4,5}$}}}{\exmath[l]{\indexedterm{\sequenceindex{3}}{i}{3}+\indexedterm{4}{i}{4}+\indexedterm{5}{i}{5}=12}}
\symrow{\displaystyle\sum_{\sequenceindex{j=2}}^{4}b_{\sequenceindex{j}}}{\conceptmeaning{Sequence terms}{\meaningexplanation{$b_{\sequenceindex{j}}=\sequenceindex{j}^2$;\\ $\sequenceindex{j=2,3,4}$}}}{\exmath[l]{\indexedterm{b_{\sequenceindex{2}}}{j}{2}+\indexedterm{b_{\sequenceindex{3}}}{j}{3}+\indexedterm{b_{\sequenceindex{4}}}{j}{4}\\[1.8mm]{}=(\sequenceindex{2})^2+(\sequenceindex{3})^2+(\sequenceindex{4})^2\\[1.8mm]{}=4+9+16=29}}
\symrow{\displaystyle\sum_{\sequenceindex{k=1}}^{4}(2\sequenceindex{k}+1)}{\conceptmeaning{Linear summand}{\meaningexplanation{$\sequenceindex{k=1,\ldots,4}$; add $2\sequenceindex{k}+1$}}}{\exmath[l]{\indexedterm{(2(\sequenceindex{1})+1)}{k}{1}+\indexedterm{(2(\sequenceindex{2})+1)}{k}{2}\\[1.8mm]{}\quad+\indexedterm{(2(\sequenceindex{3})+1)}{k}{3}+\indexedterm{(2(\sequenceindex{4})+1)}{k}{4}\\[1.8mm]{}=3+5+7+9=24}}
}
\endgroup
{\color{navy}\hrule height 0.5pt}\vspace{0.7mm}
{\FooterFont Symbols acquire meaning from their mathematical context. State the starting index and label the sequence in every partial sum.\par}
\endgroup
\newpage
\posterheader{B-2\quad LIMITS AND CALCULUS}
\begingroup
% The densest poster needs slightly smaller body/example type; symbols and
% headings retain the common sizes and no table is scaled.
\renewcommand{\DefinitionFont}{\fontsize{12.6}{14.4}\selectfont}
\renewcommand{\ExampleFont}{\fontsize{12.4}{14.4}\selectfont}
\setlength{\RowPadding}{0.20mm}
% Local column proportions reserve room for wrapped calculus examples.
\renewcommand{\SymbolColumnFraction}{.25}
\renewcommand{\MeaningColumnFraction}{.34}
% Balance body rows locally while keeping both header baselines aligned.
\newcommand{\calculusrowpadding}[1]{%
  \setlength{\RowPadding}{#1}%
  \renewcommand{\paddedcell}[1]{{\setlength{\RowPadding}{0.20mm}\centeredcell{##1}}}%
}
\sectionheading{Limits and approaching values}
\noindent\begin{minipage}[t]{.49\linewidth}\vspace{0pt}
\notationtable{
\symrow{x\mathrel{\sequenceindex{\mathord{\to}}}a}
  {\conceptmeaning{Approaching a value}{\meaningexplanation{Arbitrarily close to $a$; not $x=a$.}}}
  {$x\mathrel{\sequenceindex{\mathord{\to}}}2$: $x$ approaches $2$;\par $x=2$ is exact.}
\symrow{\begin{gathered}x\to a^{\sequenceindex{-}}\\x\to a^{\sequenceindex{+}}\end{gathered}}
  {\conceptmeaning{One-sided approaches}{\meaningexplanation{Inputs below ($\sequenceindex{\mathord{-}}$) or above ($\sequenceindex{\mathord{+}}$) $a$; not signs of $f(x)$.}}}
  {\exmath[l]{x\to2^{\sequenceindex{-}}:\ x<2;\\x\to2^{\sequenceindex{+}}:\ x>2.}}
\symrow{\begin{gathered}x\to\sequenceindex{\mathord{+}\infty}\\[1pt]x\to\sequenceindex{\mathord{-}\infty}\end{gathered}}
  {\conceptmeaning{Special cases: inputs without bound}{\meaningexplanation{Inputs unbounded in either direction.}}}
  {\exmath[l]{\lim\limits_{x\to\sequenceindex{\mathord{+}\infty}}\frac1x=0,\\[1pt]\lim\limits_{x\to\sequenceindex{\mathord{-}\infty}}\frac1x=0.}}
\symrow{\begin{gathered}\mathop{\sequenceindex{\mathord{\lim}}}\limits_{x\to a}f(x)\\{}=L\end{gathered}}
  {\conceptmeaning{Two-sided finite limit}{\meaningexplanation{As domain inputs $x\ne a$ approach $a$, both sides tend to finite $L$.}}}
  {\exmath[l]{\mathop{\sequenceindex{\mathord{\lim}}}\limits_{x\to2^-}x^2=4\\{}=\mathop{\sequenceindex{\mathord{\lim}}}\limits_{x\to2^+}x^2\\[1pt]
   \mathop{\sequenceindex{\mathord{\lim}}}\limits_{x\to2}x^2=4;\\[1pt]\mathop{\sequenceindex{\mathord{\lim}}}\limits_{x\to1}\frac{x^2-1}{x-1}=2}
   \par The quotient is undefined at $x=1$.}
}
\end{minipage}\hfill
\begin{minipage}[t]{.49\linewidth}\vspace{0pt}
\renewcommand{\SymbolColumnFraction}{.32}
\renewcommand{\MeaningColumnFraction}{.32}
\calculusrowpadding{1.1mm}
\notationtable{
\symrow{\begin{gathered}f(x)\to\sequenceindex{\mathord{+}\infty}\\[1pt]f(x)\to\sequenceindex{\mathord{-}\infty}\end{gathered}}
  {\conceptmeaning{Special case: unbounded values}{\meaningexplanation{Unbounded values; $\infty$ is not a finite real limit.}}}
  {\exmath[l]{\lim\limits_{x\to a}\frac1{(x-a)^2}=\sequenceindex{\mathord{+}\infty}\\[1pt]
   \lim\limits_{x\to0^+}\frac1x=\sequenceindex{\mathord{+}\infty};\\[1pt]\lim\limits_{x\to0^-}\frac1x=\sequenceindex{\mathord{-}\infty}}}
\symrow{\begin{gathered}\lim\limits_{x\to a}f(x)\\\sequenceindex{\text{DNE}}\end{gathered}}
  {\conceptmeaning{Special case: limit does not exist (DNE)}{\meaningexplanation{Unequal one-sided limits. Here $g(x)=|x|/x$, $x\ne0$.}}}
  {\exmath[l]{\lim\limits_{x\to0^-}g(x)=-1\\{}\ne1=\lim\limits_{x\to0^+}g(x)\\[1pt]
   \lim\limits_{x\to0}g(x)\ \sequenceindex{\text{DNE}}}}
\symrow{\begin{gathered}f^{\sequenceindex{\prime}}(a)=\\\displaystyle\lim\limits_{h\to0}\frac{f(a+h)-f(a)}h\end{gathered}}
  {\conceptmeaning{Definition of derivative}{\meaningexplanation{Requires a finite limit; $h\ne0$.}}}
  {For $f(x)=x^2$:\par
   \exmath[l]{f^{\sequenceindex{\prime}}(a)=\lim\limits_{h\to0}(2a+h)\\{}=2a.}}
}
\end{minipage}\par\vspace{2mm}
\sectionheading{Change and derivatives}
\renewcommand{\ContextHeading}{A derivative is a limit}
\examplecontexts{For $y=f(x)$, the derivative gives instantaneous change. Its limit definition appears above; the rows below distinguish finite changes, differentiation, and derivative values.}
\noindent\begin{minipage}[t]{.49\linewidth}\vspace{0pt}
\notationtable{
\symrow{\begin{gathered}\sequenceindex{\Delta} x,\\\sequenceindex{\Delta} y\end{gathered}}
  {\conceptmeaning{Finite changes}{\meaningexplanation{$\sequenceindex{\Delta} x=x_2-x_1$; $\sequenceindex{\Delta} y=f(x_2)-f(x_1)$.}}}
  {$f(x)=x^2$, from $x=1$ to $x=3$:\par $\sequenceindex{\Delta} x=2$, $\sequenceindex{\Delta} y=9-1=8$.}
\symrow{\frac{\sequenceindex{\Delta} y}{\sequenceindex{\Delta} x}}
  {\conceptmeaning{Average rate of change}{\meaningexplanation{Output change per input change; $\sequenceindex{\Delta} x\ne0$.}}}
  {$f(x)=x^2$, from $1$ to $3$:\par $\sequenceindex{\Delta} y/\sequenceindex{\Delta} x=8/2=4$.}
\symrow{\begin{gathered}f^{\sequenceindex{\prime}}(x),\quad y^{\sequenceindex{\prime}},\\\frac{\sequenceindex{d}y}{\sequenceindex{d}x}\end{gathered}}
  {\conceptmeaning{First derivative}{\meaningexplanation{Instantaneous rate for $y=f(x)$.}}}
  {$y=f(x)=x^2$:\par $\begin{gathered}f^{\sequenceindex{\prime}}(x)=y^{\sequenceindex{\prime}}\\{}=\dfrac{\sequenceindex{d}y}{\sequenceindex{d}x}=2x\end{gathered}$.}
}
\end{minipage}\hfill
\begin{minipage}[t]{.49\linewidth}\vspace{0pt}
\calculusrowpadding{1mm}
\notationtable{
\symrow{\frac{\sequenceindex{d}}{\sequenceindex{d}x},\quad f^{\sequenceindex{\prime}}(a)}
  {\conceptmeaning{Differentiation; value at a point}{\meaningexplanation{$\sequenceindex{d}/\sequenceindex{d}x$: operator; $f^{\sequenceindex{\prime}}(a)$: derivative value at $a$.}}}
  {$\dfrac{\sequenceindex{d}}{\sequenceindex{d}x}(x^2)=2x$;\par $f^{\sequenceindex{\prime}}(3)=2(3)=6$.}
\symrow{\begin{gathered}f^{\sequenceindex{\prime\prime}}(x),\\\frac{d^{\sequenceindex{2}}y}{dx^{\sequenceindex{2}}}\end{gathered}}
  {\conceptmeaning{Second derivative}{\meaningexplanation{Derivative of the first derivative.}}}
  {$y=f(x)=x^2$:\par $\begin{gathered}f^{\sequenceindex{\prime\prime}}(x)\\{}=\dfrac{d^{\sequenceindex{2}}y}{dx^{\sequenceindex{2}}}=2\end{gathered}$.}
}
\end{minipage}\par\vspace{2mm}
\sectionheading{Integrals and evaluation}
\noindent\begin{minipage}[t]{.49\linewidth}\vspace{0pt}
\notationtable{
\symrow{\notationintegral{\int} f(x)\,dx}
  {\conceptmeaning{Indefinite integral}{\meaningexplanation{Antiderivatives $F(x)+C$; $F'=f$, $C\in\mathbb R$.}}}
  {$\displaystyle\notationintegral{\int} 2x\,dx=x^2+C$\par A family of functions, one for each $C$.}
\symrow{\notationintegral{\int_a^b} f(x)\,dx}
  {\conceptmeaning{Definite integral}{\meaningexplanation{Signed accumulation from $a$ to $b$; integrand $f(x)$.}}}
  {$\displaystyle\notationintegral{\int_0^2} 2x\,dx=4$\par A numerical value; signed accumulation is not always geometric area.}
\symrow{F(x)\sequenceindex{\big|}_a^b}
  {\conceptmeaning{Evaluation at the bounds}{\meaningexplanation{$F(b)-F(a)$. Equals $\int_a^b f(x)\,dx$ if $f$ is continuous on $[a,b]$ and $F'=f$.}}}
  {\exmath[l]{\displaystyle\int_0^2 2x\,dx=x^2\sequenceindex{\big|}_0^2\\{}=2^2-0^2=4}.}
}
\end{minipage}\hfill
\begin{minipage}[t]{.49\linewidth}\vspace{0pt}
\calculusrowpadding{1mm}
\notationtable{
\symrow{\begin{gathered}\sequenceindex{\Delta x}=\frac{b-a}{n}\\
  \sequenceindex{x_i}=a+i\sequenceindex{\Delta x}\\
  \sequenceindex{R_n}=\\\sum_{i=1}^n f(\sequenceindex{x_i})\sequenceindex{\Delta x}\end{gathered}}
  {\conceptmeaning{Riemann sum}{\meaningexplanation{Continuous $f$ on $[a,b]$; $n$ equal widths $\sequenceindex{\Delta x}$, right endpoints $\sequenceindex{x_i}$.}\par
  $\begin{gathered}\displaystyle\int_a^b f(x)\,dx\\{}=\lim\limits_{n\to\infty}\sequenceindex{R_n}\end{gathered}$.}}
  {For $f(x)=2x$ on $[0,2]$:\par
   \exmath[l]{\sequenceindex{\Delta x}=\sequenceindex{2/n},\\\sequenceindex{x_i}=\sequenceindex{2i/n}\\
   \sequenceindex{R_n}=\sum_{i=1}^n2\bigl(\sequenceindex{2i/n}\bigr)\bigl(\sequenceindex{2/n}\bigr)\\
   {}=\sequenceindex{4(n+1)/n}\to\sequenceindex{4}}}
\symrow{\sequenceindex{d}x}
  {\conceptmeaning{Integration variable}{\meaningexplanation{Integrate with respect to $x$.}}}
  {In $\int_0^2 2x\,\sequenceindex{d}x$: bounds $0$, $2$;\par integrand $2x$; integration variable $x$.}
}
\end{minipage}\par

\par\vspace{2mm}
\sectionheading{Relevant examples}
\begingroup
\setlength{\tabcolsep}{1mm}
\setlength{\RowPadding}{0.45mm}
\renewcommand{\arraystretch}{1}
\arrayrulecolor{rulecolor}
\newcommand{\calculusrow}[3]{%
  \paddedcell{\SymbolFont $\displaystyle #1$}&%
  \paddedcell{\SymbolFont $\displaystyle #2$}&%
  \paddedcell{\SymbolFont $\displaystyle #3$}\\\hline}
\noindent\begin{minipage}[t]{.52\linewidth}\vspace{0pt}
\begin{tabularx}{\linewidth}{@{}p{.30\linewidth}p{.26\linewidth}X@{}}
\paddedcell{\TableHeadFont Function $f(x)$}&
\paddedcell{\TableHeadFont Derivative $f'(x)$}&
\paddedcell{\TableHeadFont Indefinite integral\par $\int f(x)\,dx$}\\\hline
\calculusrow{e^x}{e^x}{e^x+C}
\calculusrow{\ln x\quad(x>0)}{\frac1x}{x\ln x-x+C}
\calculusrow{P(x)=\sum_{k=0}^{n}a_kx^k}
  {\sum_{k=1}^{n}k a_kx^{k-1}}
  {\sum_{k=0}^{n}\frac{a_k}{k+1}x^{k+1}+C}
\end{tabularx}
\end{minipage}\hfill
\begin{minipage}[t]{.46\linewidth}\vspace{0pt}
\begin{tabularx}{\linewidth}{@{}p{.28\linewidth}p{.29\linewidth}X@{}}
\paddedcell{\TableHeadFont Function $f(x)$\par $x$ in radians}&
\paddedcell{\TableHeadFont Derivative $f'(x)$}&
\paddedcell{\TableHeadFont Indefinite integral\par $\int f(x)\,dx$}\\\hline
\calculusrow{\sin x}{\cos x}{-\cos x+C}
\calculusrow{\cos x}{-\sin x}{\sin x+C}
\calculusrow{\begin{gathered}\tan x\\(\cos x\ne0)\end{gathered}}{\begin{gathered}\sec^2 x=\\\frac1{\cos^2 x}\end{gathered}}{-\ln|\cos x|+C}
\end{tabularx}
\end{minipage}\par\vspace{0.8mm}
\endgroup
{\color{navy}\hrule height 0.5pt}\vspace{0.7mm}
{\FooterFont Limits describe approaching behavior; derivatives describe local change; definite integrals describe signed accumulation.\par}
\endgroup
\end{document}
